% 主文件。导入 ZIP 后可直接 Recompile，无需手动更改编译器。
\documentclass{gmcmthesis}

% 题目：三号黑体。加长下划线后，文字仍居中，因而向右移动。
\title{论文题目}
\setlength{\papertitleleft}{148.1bp} % 下划线左端；增大则整条线和文字一起右移
\setlength{\papertitlewidth}{300bp}  % 下划线长度；原模板为245.03bp

% 封面填写区；留空时与附件3的空白封面保持一致。
\schoolname{}
\baominghao{}
\membera{}
\memberb{}
\memberc{}

\begin{document}
\makecover
\maketitle
\begin{abstract}
  在此填写摘要。开头概括问题背景、研究目标与总体建模思路，再按赛题顺序说明各问题采用的方法、模型及求解过程，最后给出有依据的主要结果、结论与创新点。

针对问题一，在此填写核心假设、模型构建方法与实际计算结果。摘要中的数据应与正文、图表和程序输出一致。

针对问题二，在此说明模型如何处理新增约束，并概括求解方法及验证结果。按实际赛题增减段落，不保留这些写作提示。

最后概括模型适用范围和检验结论。摘要一般不超过两页，无需另附英文摘要。本页文字仅演示字体、段落和缩进，没有提供任何研究结果。

\keywords{关键词一\quad 关键词二\quad 关键词三}

\end{abstract}

% 附件2要求摘要后下一页开始正文，因此默认不插入目录。
% 若提交方另有目录要求，可在此启用 \maketoc。
\startbody
\section{问题重述}
\subsection{问题背景}
在此说明赛题所研究对象、应用场景和需要解决的实际问题。保留与模型建立直接有关的信息，用本队语言概括背景。

\subsection{问题描述}
\textbf{问题一：}在此概括第一个问题的目标、输入和预期输出。

\textbf{问题二：}在此概括第二个问题的目标、约束及其与问题一的联系。

\textbf{问题三：}根据实际赛题填写，可增删问题数量。

\section{问题分析}
\subsection{总体思路}
在此说明各问题之间的关系，解释为什么选择相应的数学方法。可用流程图展示数据处理、模型建立、求解、验证之间的顺序。

\subsection{各问题的关键难点}
按问题说明建模难点、可用数据和应对策略，使后续模型的选择有明确依据。

\section{模型假设}
\begin{enumerate}[label=\arabic*.]
  \item 在此填写必要的模型假设，并说明其适用条件。
  \item 在此填写下一项假设，避免把需要检验的结论写成前提。
\end{enumerate}

\section{符号说明}
\begin{table}[htbp]
  \centering
  \caption{主要符号及含义}
  \label{tab:notations}
  \begin{tabularx}{\textwidth}{cXc}
    \toprule
    符号 & 含义 & 单位 \\
    \midrule
    $n$ & 样本数量（演示符号，请替换） & 无量纲 \\
    $x_i$ & 第 $i$ 个观测值（演示符号，请替换） & 按实际填写 \\
    $\theta$ & 模型参数（演示符号，请替换） & 按实际填写 \\
    \bottomrule
  \end{tabularx}
\end{table}
仅列出全文反复使用的重要符号，局部符号可在首次出现处说明。

\section{模型的建立与求解}
\subsection{数据预处理}
说明数据来源、清洗规则、缺失值与异常值处理，以及后续计算使用的数据形式。数据统计以实际处理结果为准。

\subsection{问题一的模型建立与求解}
\subsubsection{模型建立}
在此给出变量、目标函数、约束条件和推导过程。下面的均值公式仅用于演示公式排版，并非针对赛题提出的模型：
\begin{equation}
  \bar{x}=\frac{1}{n}\sum_{i=1}^{n}x_i .
  \label{eq:example}
\end{equation}
公式\eqref{eq:example}中的符号应在正文或符号表中解释。

\subsubsection{求解方法}
在此说明算法步骤、参数设置、停止条件与计算环境。必要时说明数值稳定性、复杂度或可复现的方法。

\subsubsection{结果与分析}
在此填写真实求解结果，并通过图表进行分析。下表只有占位文字，不代表实验结果。
\begin{table}[htbp]
  \centering
  \caption{结果展示格式示例}
  \label{tab:results}
  \begin{tabularx}{\textwidth}{lXXX}
    \toprule
    方案 & 指标一 & 指标二 & 说明 \\
    \midrule
    方案一 & 待计算 & 待计算 & 待填写 \\
    方案二 & 待计算 & 待计算 & 待填写 \\
    \bottomrule
  \end{tabularx}
\end{table}

\begin{figure}[htbp]
  \centering
  \fbox{\parbox[c][35mm][c]{0.8\textwidth}{\centering 在此插入实际结果图}}
  \caption{结果图排版示例}
  \label{fig:example}
\end{figure}
图\ref{fig:example}的图注位于图下方，表\ref{tab:results}的表题位于表上方。

\subsection{问题二的模型建立与求解}
\subsubsection{模型建立}
在此说明第二个问题的新增变量、约束及模型。
\subsubsection{求解与结果分析}
在此填写求解方法、真实结果与解释。

\subsection{问题三的模型建立与求解}
按实际问题补充模型、求解过程与结果。赛题包含更多问题时，保持同样的标题层级。

\section{模型检验与灵敏度分析}
\subsection{模型检验}
根据任务采用合理的误差检验、交叉验证、基准比较或约束检查。列出真实结果并解释其意义。
\subsection{灵敏度与稳健性分析}
说明关键参数、输入扰动或假设变化对结果的影响，给出模型适用范围。

\section{模型评价与推广}
\subsection{模型优点}
用分析或结果支持模型的优点，避免没有证据的评价。
\subsection{模型局限与改进}
说明数据、假设和方法的局限，以及可行的改进方向。

\section{结论}
在此逐项回应赛题目标，总结已由结果支持的结论。不要引入正文没有论证的新结论。

% 按实际引用顺序填写真实来源；以下是格式占位，不是文献。
% 正文使用 \cite{ref:actual} 引用，下方替换该条目。
\begin{thebibliography}{99}
  \bibitem{ref:actual}
  作者，真实文献题名，出版信息或网址，页码或访问日期（请替换为实际引用来源）。
\end{thebibliography}

\appendicesbegin
\section{补充材料与程序}
在此列出必要的补充推导、程序清单和运行说明。正文引用本附录时使用交叉引用。
\subsection{程序排版示例}
以下代码仅演示程序格式，应替换为实际求解程序。
\begin{lstlisting}[language=Python]
def main():
    pass

if __name__ == "__main__":
    main()
\end{lstlisting}

\end{document}
